\section{Introduction}
In this section, we discuss common concepts that are used frequently in reinforcement learning.

\section{State, Action, Policy and Reward}
\subsection{State}
    The \textbf{state} is the status of the agent with respect to the environment.
    The set of all possible states, the \textbf{state space}, is denoted as $\mathcal{S}$.
\subsection{Action}
    The \textbf{action} is what agent can do, when a state is given. By choosing the action and
    taking it, the agent moves to the next state, which is called \textbf{state transition}. The set of all possible actions in state $s$,
    \textbf{action space} for state $s$, is denoted as $\mathcal{A}(s)$.
    \subsubsection{Transition Probability}
        Note that the transition can be stochastic. That is, the same choice of action in the same state
        may not lead to equivalent next state. We denote this probability with conditional porbability as following:
        \begin{equation}
            p(s'|s,a)
        \end{equation}
        where $s',s \in \mathcal{S}$ and $a \in \mathcal{A}(s)$.
\subsection{Policy}
    The \textbf{policy} tells the agent which actions to take at every state. The policy can be deterministic or stochastic.
    We denote the probability of taking action $a$ under state $s$ following policy $\pi$ as $\pi(a|s)$.
\subsection{Reward}
    \textbf{Reward} is one of the most unique concept of reinforcement learning. After the execution of action, the agent
    receives a reward, denoted as $r$. The agent's goal is to maximize the total reward that the agent receives.
    Due to this goal, the reward can encourage or discourage the agent's
    choice of selecting an action under certain state. Again, the reward can be stochastic, too.

    One question that follows naturally is: Can't we just pick the actions that gives the most reward at every state? If this was correct,
    the reinforcement learning may not have existed. Since the reward is just an immediate information of the act, choosing only the action with
    highest reward is not the correct answer, most of the time. There can be a better choice in the long run.

\section{Trajectories, Returns and Episodes}
    \subsection{Trajectory}
        The \textbf{trajectory} is the footsteps of the agent. More formally, suppose that the agent is on the initial state $s_0$ and took action $a_0$.
        The agent receives reward $r_1$ from environment and moves to state $s_1$. Again, the agent take an action $a_1$, receive reward $r_1$ and move to state $s_2$.
        Repeating this, we get the trajectory of an agent:
        \begin{equation}
            s_0,a_0,r_1,s_1,a_1,r_2,\dots,s_{T-1},a_{T-1},r{T}. \label{eq:trajectory}
        \end{equation}
    \subsection{Returns}
        Given the trajectory $\tau = \{s_0,a_0,r_1,s_1,a_1,r_2,\dots,s_{T-1},a_{T-1},r{T}\}$, we can calculate the \textbf{return} $G_t$, of the trajectory.
        The naive definition of the return, called \textbf{total rewards} is the sum
        \begin{equation}
            G_t = \sum_{k=t+1}^{T}r_k. \label{eq:total-reward}
        \end{equation}
        The point is that we do not know if the terminal time $T$ is finite or not. If $T$ is not finite, than the cummulative reward \ref{eq:total-reward}
        diverges. For this, we introduce the \textbf{discounted return}, defined as
        \begin{equation}
            G_t = \sum_{k=t+1}^{\infty}\gamma^{k - t - 1}r_k
        \end{equation}
        where $\gamma \in [0,1]$ is called \textbf{discount rate}.
        This works for for finite trajectory too, if we let $r_k = 0$ for $k > T$.
    \subsection{Episodes}
        The \textbf{episode} is the finite trajectory itself, like \ref{eq:trajectory}. The episode may stop at the states called \textbf{terminal states}.
        The tasks with episodes are called \textbf{episodic tasks}. Unlike episodic tasks, there are \textbf{continuing tasks}, which does not ends. These two can
        be expressed in unified mathematical manner. For this, in the episodic task, we must consider the agent's behavior after the terminal state:
        to stay in that state forever, or to take action like normal states. In this note, we follow the latter, like \cite{Zhao2025}.

\section{Markov Decision Process}
    \subsection{Definition}
        The \textbf{Markov decision process} (MDP) is the general framework that describes the stochastic dynamic systems. It is defined as following:
        \begin{defn}{Markov Decision Process}{MDP}
            The \textbf{Markov decision process} is the 5-tuple $(\mathcal{S},\mathcal{A},p,\mathcal{R},\gamma)$, where each elements are defined as following:
            \begin{itemize}
                \item $\mathcal{S}$ is the state space, $\mathcal{A}(s)$ is the action space associated with state $s$, and $\mathcal{R}(s,a)$ is a set of rewards, associated with each state-action pair $(s,a)$.
                \item $p(s',r|s,a)$ is the dynamics of the model, which is as probability of moving to state $s'$ from $s$ by taking action $a$, and obtaining reward $r$.
                Note that the state transition probability and the reward probability can be decribed with dynamics. See \ref{rem:dynamics}
                \item $\pi(s|a)$ is the policy, a probability of agent choosing action $a$ in state $s$.
                \item $\gamma$ is the discount rate.
            \end{itemize}
            with the dynamics satisfying the \emph{Markov property}:
            \begin{align}
                p(s_{t+1}|s_t,a_t,s_{t-1},a_{t-1},\dots,s_0,a_0) &= p(s_{t+1}|s_{t},a_{t}) \\
                p(r_{t+1}|s_t,a_t,s_{t-1},a_{t-1},\dots,s_0,a_0) &= p(r_{t+1}|s_{t},a_{t})
            \end{align}
            where $t$ is the current time step.
        \end{defn}
        \begin{rem}{Transition and reward probability as dynamics}{dynamics}
            The 4-tuple dynamics can be used to describe transition probability and reward probability as following:
            \begin{align}
                p(s'|s,a) &= \sum_{r \in \mathcal{R}(s,a)}p(s',r|s,a) \\
                p(r|s,a) &= \sum_{s' \in \mathcal{S}}p(s',r|s,a)
            \end{align}
        \end{rem}
        The dynamics can be either stationary or non-stationary. If the dynamics changes over time, it is non-stationary; else, it is called stationary.
